\begin{prob}
    Solve the following recurrence relations.
    
    \begin{subprobset}
    
            \begin{subprob}
		        T(n) = T(n-1) + n
		        \\
		        T(0)=0
                \begin{soln}
                    \begin{align*}
                        T(n) &= T(n-1) + n
                        \\
                        &= [T(n-2) + n-1] + n
                        \\
                        &= T(n-2) + 2n - 1
                        \\	
                        &= [T(n-3) + n-2] + 2n - 1
                        \\
                        &= T(n-3) + 3n - (1 + 2)
                        \\
                        &= [T(n-4) + n -3] + 3n - (1+2)
                        \\
                        &= T(n-4) + 4n - (1+2+3)
                    \end{align*}
                    We can infer that $T(n) = T(n - k) + kn - \sum_{i=1}^{k-1}i$ in the k-th step. 
                    \\
                    T(0) is the base case. $n-k =0 $ when $n=k$.
                    \\
                    \begin{align*}
                        \sum_{i=1}^{n-1}i = \frac{(n-1)n}{2} = \frac{n^2 - n}{2}
                    \end{align*}
                    \begin{align*}
                        T(n) &= T(n-n) + n\cdot n - \sum_{i=1}^{n-1}i 
                        \\
                        &= T(0) + n^2 -  \frac{n^2 - n}{2}
                        \\
                        &= 0 + n^2 -  \frac{n^2 - n}{2}
                        \\  
                        &= \theta(n^2)
                    \end{align*} 
                \end{soln}
            \end{subprob}
            
            \vspace{25em}
            \begin{subprob}
                T(n)=4T(n/4) + n
                \\
                T(1)=1
                \begin{soln}
                    \begin{align*}
                        T(n) &= 4 \cdot T(n/4) + n
                        \\
                        &=
                        4 \left[
                        4 \cdot T(n/16) + n/4
                        \right] + n
                        \\
                        &=
                        16 \cdot T(n/16) + 2n
                        \\
                        &=
                        16 \left[
                        4 \cdot T(n/64) + n/16
                        \right] + 2n
                        \\
                        &=
                        64 \cdot T(n/64) + 3n
                        \intertext{We can infer that in the k-th step, we have:}
                        &=
                        4^k \cdot T(n/4^k) + k\cdot n
                    \end{align*}
                    The base case will be reached when $n/4^k = 1$, that is, when $k = \log_4 n$.
                    Substituting this value of $k$ into the general expression:
                    \begin{align*}
                        T(n)
                        &=
                        4^{\log_4 n} \cdot T(n/4^{\log_4 n}) + n \cdot \log_4 n
                        \\
                        &=
                        n \cdot T(n/n) + n \cdot \log_4 n
                        \\
                        &=
                        n \cdot T(1) + n \cdot \log_4 n
                        \intertext{Since $T(1) = 1$, we have:}
                        &=
                        n + n \cdot \log_4 n\\
                        &=
                        \Theta(n \log_4 n)\\
                        \intertext{Since logarithms of different bases differ only by a constant
                        factor, we typically omit the base when using asymptotic notation:}
                        &=
                        \Theta(n \log n)
                    \end{align*}
                \end{soln}
        \end{subprob}

    \end{subprobset}

\end{prob}
